\subsection{局限性与未来工作}
\label{sec:limitations}

本文仍存在若干限制。首先，实验场景为静态合成低空环境，主要用于验证强约束三维路径规划中的可行性形成、轨迹质量和评价开销，尚未引入真实地图、DEM 数据或城市级障碍物建模。其次，当前任务假设障碍物和禁飞区在规划周期内保持不变，尚未覆盖动态障碍、在线重规划或多无人机协同避碰。再次，当前实现对不可行解采用实现相关违反分数排序，并使用保守局部 CVF 候选接收规则；归一化违反量、安全优先词典序或更激进的可行性改善接收均属于后续算法变体，需要重新开展正式实验。最后，本文评估对象为路径规划层面的几何可行性和优化性能，尚未与飞控约束、动力学模型和轨迹跟踪控制器进行闭环联调。后续工作可在真实地形数据、动态低空交通约束和控制器可跟踪性验证中进一步检验 CVF-AE 的工程适用性。
